\section{Related Work}\label{sec:related}

\noindent\textbf{Monte Carlo Localization.}
Particle-filter localization was introduced by \citet{dellaert1999montecarlo} and \citet{fox1999efficient}, who represent the pose posterior as a weighted sample set updated by motion-model prediction, sensor-likelihood weighting, and resampling, and who demonstrated global localization from a uniform prior on real robots. \citet{thrun2001robust} later hardened MCL against accurate sensors and the kidnapped-robot problem via Mixture-MCL's dual sampler, and \citet{fox2003kld} made the sample size adaptive through KLD-sampling; the odometry motion model and likelihood-field sensor model we use are the textbook forms of \citet{thrun2005probabilistic}. The particle-filter core of \sname{} is a direct descendant of this line and of its \code{nav2\_amcl} lineage; we contribute no new filter theory. What differs is factoring: the same predict--weight--resample core is shared, unchanged, between a 2D likelihood-field backend and a 3D NDT backend, and recovery is delegated to a deterministic branch-and-bound search (\Cref{sec:bbs}) rather than stochastic injection or dual sampling.

\noindent\textbf{NDT Representations and Measurement Models.}
\citet{biber2003ndt} introduced the Normal Distributions Transform, modeling a scan as per-cell Gaussians and reducing registration to Newton optimization of a smooth score; \citet{magnusson2007scan} generalized the representation to 3D voxel grids for field deployment on mining vehicles. Both use NDT for deterministic scan-to-scan or scan-to-map registration that requires a reasonable initial guess. The \code{ndt3d} backend of \sname{} adopts the Gaussian-voxel representation but reads the NDT score out as a per-particle likelihood inside the shared MCL core (NDT-MCL rather than NDT registration), trading the quadratic convergence of Newton alignment for the multi-hypothesis robustness of the particle filter, so the 2D and 3D backends differ only in their measurement model.

\noindent\textbf{LiDAR-Inertial Fusion and Robot Software Systems.}
\citet{sola2017quaternion} is the mathematical basis of the \code{fusion3d} backend: the nominal/error-state decomposition, local angular-error parameterization, and predict--correct--inject--reset cycle are implemented essentially as derived there. On the systems side, \citet{moore2014generalized} established the configurable EKF/UKF fusion node for ROS, \citet{wan2018robust} demonstrated an industrial ESKF fusing GNSS-RTK, LiDAR, and IMU at centimeter accuracy as a production system tied to its own pre-built LiDAR maps, and Nav2 \citep{macenski2020marathon} treats localization as an external, swappable component. Tightly-coupled LiDAR-inertial odometry such as FAST-LIO \citep{xu2021fastlio}, and applied 3D map-based tracking systems \citep{koide2019portable}, are complementary: they produce ego-motion and maps, whereas \sname{} localizes within a prior map. Relative to these, \sname{} places three known estimator families behind one ROS~2 topic and TF contract, each implemented as a library that is unit-testable without a middleware runtime, a property that \code{robot\_localization}'s otherwise modular design did not prioritize, and which answers the modular-localization framework that \citet{macenski2023survey} left as future work.

\noindent\textbf{Global (Kidnapped-Robot) Localization.}
\citet{olson2009correlative} formulated 2D scan matching as exhaustive correlative search over a rasterized log-likelihood table with coarse-to-fine multi-resolution pruning, and \citet{hess2016realtime} bounded that search with branch-and-bound over a precomputed grid pyramid to run loop closure in real time inside Cartographer. \sname{} borrows exactly this branch-and-bound-over-pyramid primitive but applies it in a narrower setting: relocalization against a known static map, with no submaps or pose graph to maintain, whose verified result re-seeds the particle filter in place of the heuristic recovery of classical MCL \citep{fox1999efficient,thrun2001robust}.
